\documentclass[10pt, a4paper, twocolumn]{article} % 两栏排版，最省纸

% --- 基础设置（这部分你永远不需要动） ---
\usepackage[UTF8]{ctex} % 中文支持
\usepackage[left=1cm, right=1cm, top=1.5cm, bottom=1.5cm]{geometry} % 极窄边距
\usepackage{listings} % 代码显示
\usepackage{xcolor}
\lstset{
    language=C++,
    basicstyle=\ttfamily\scriptsize, % 字体极小，为了省纸
    breaklines=true,
    numbers=left,
    frame=single,
    tabsize=2
}

\begin{document}

\title{我的竞赛参考资料 (TRD)}
\author{姓名}
\maketitle
\tableofcontents % 自动生成目录
\newpage

% --- 【这里开始是你的填空区】 ---

\section{杂项}
    \subsection{trick}
    \subsubsection{区间排序+扫描}
    \begin{lstlisting}
        将所有区间按照左端点升序排序，然后维护最大右端点的套路在很多题都会用到。注意要维护前缀的最大右端点，不能只是拿当前右端点跟上一个区间的右端点比较，因为右端点不一定是升序的。
        然后对于上面的这种套路右端点升序还是降序通常不影响(就是说如果左端点相同那么右端点升序还是降序排列)，但是在一些其他的区间问题中可能次关键字(这里的右端点)的排序方向很重要，不能一概而论。
    \end{lstlisting}
    \subsection{板子}
    \subsubsection{最小质因子筛法}    
    \begin{lstlisting}
        const int MAXN=200005;
		int min_prime[MAXN];//min_prime[i]表示i的最小质因数
		void sieve()//预处理所有数的最小质因数
		{
			for(int i=2;i<MAXN;i++)
			{
				if(min_prime[i]==0)//表示i是质数，因为如果i是合数，那么它肯定在之前已经被标记过了   
				{
					for(int j=i;j<MAXN;j+=i)//把i的倍数还没被标记的全部标记上
					{
						if(min_prime[j]==0) min_prime[j]=i;
					}
				}
			}
		}
		vector<int>get_prime_factors(int n)//快速获取一个数的所有不同质因数，用vector形式返回
		{
			vector<int>factors;
			while(n>1)
			{
				int p=min_prime[n];//找当前n的最小质因数
				factors.pb(p);
				while(n%p==0) n/=p;//去掉所有相同的质因子
			}
			return factors;
		}
    \end{lstlisting}

    \subsubsection{坐标离散化}
    \begin{lstlisting}
    	//坐标压缩：数组的值可能很大，达到1e9，这时候用vector会爆，用map常数太大，可以用坐标压缩。
           //如果我们不在意数组的具体值，只在乎有多少个不同的数值，那么让每个数值=他在数组中是第几大的数，这样构成的新数组最大也就2e5。
            vector<int>a(n);
			vector<int>b;
			for(int i=0;i<n;i++)
			{
				cin>>a[i];
				b.pb(a[i]);
			}
			sort(b.begin(),b.end());
			b.erase(unique(b.begin(),b.end()),b.end());
			for(int i=0;i<n;i++)
			{
				a[i]=lower_bound(b.begin(),b.end(),a[i])-b.begin();
			}
    \end{lstlisting}


    \subsubsection{exgcd}
    \begin{lstlisting}
int exgcd(int a,int b,int &x,int &y)//得到的x和y是右边为gcd（a，b）时的特解，如果右边是k*gcd（a，b），还要将x，y乘k
{
    if(b==0)
    {
        x=1,y=0;
        return a;
    }
    int d=exgcd(b,a%b,y,x);
    y-=(a/b)*x;
    return d;
}
//通解形式可以表示为：x=xp+k*（b/g）
//                  y=yp+k*(a/g)
    \end{lstlisting}
    \subsubsection{随机数}
    \begin{lstlisting}
mt19937_64 rng(chrono::steady_clock::now().time_since_epoch().count());//在全局写这句
int x=rng()%n;//生成一个0~n-1的随机数

vector<int>v;
shuffle(v.begin(),v.end(),rng);//将v中的排列随机打乱
    \end{lstlisting}
    \subsubsection{线性基}
    \begin{lstlisting}
    // 全局定义基数组
ll p[62]; 
bool has_zero = false; // 是否能异或出 0

// 1. 插入：这是最核心的
void insert(ll x) {
    for (int i = 62; i >= 0; i--) {
        if (!(x >> i & 1)) continue;
        if (!p[i]) {
            p[i] = x;
            return;
        }
        x ^= p[i];
    }
    has_zero = true; 
}

// 2. 查询最大异或和
ll query_max() {
    ll res = 0;
    for (int i = 62; i >= 0; i--) {
        if ((res ^ p[i]) > res) res ^= p[i];
    }
    return res;
}

// 3. 检查某个数能否被凑出来
bool check(ll x) {
    for (int i = 62; i >= 0; i--) {
        if (!(x >> i & 1)) continue;
        if (!p[i]) return false;
        x ^= p[i];
    }
    return true;
}

//如果想知道每一位是由idx的哪些下标组成的可以这样写，mask是定义的全局数组（用二进制形式记录每一位由哪些下标组成），local_id是idx里的下标，idx是事先判断出来的线性基的原数组的下标数组，idx最多30个元素（一般）
int mask[32];
void insert_mask(int val,int local_id)
{
	int cur_mask=(1ll<<local_id);
	for(int i=30;i>=0;i--)
	{
		if(!(val>>i&1)) continue;
		if(!p[i])
		{
			p[i]=val;
			mask[i]=cur_mask;
			return;
		}
		val^=p[i];
		cur_mask^=mask[i];
	}
}
    \end{lstlisting}
    \subsubsection{组合数}
    \begin{lstlisting}
    //帕斯卡三角形法，适用于N较小（如N<=5000）且没有模数或模数不是质数的情况。
int C[505][505]={0};//这里自己灵活改数组大小，如果没有mod太大了会爆
void init_pascal()
{
    for(int i=0;i<=500;i++)
    {
        C[i][0]=1;
        for(int j=1;j<=i;j++)
        {
            C[i][j]=(C[i-1][j-1]+C[i-1][j])%MOD;
        }
    }
}

//直接套公式法，更常用，适用于N比较大且有模数并且模数是质数的情况
int fact[MAXN], invFact[MAXN];


int fastPow(int a, int b) {
    int res = 1;
    a %= MOD;
    while (b > 0) {
        if (b % 2 == 1) res = (res * a) % MOD;
        a = (a * a) % MOD;
        b /= 2;
    }
    return res;
}


int modInverse(int n) {
    return fastPow(n, MOD - 2);
}

// 预处理阶乘和阶乘逆元
void precompute() {
    fact[0] = 1;
    for (int i = 1; i < MAXN; i++) {
        fact[i] = (fact[i - 1] * i) % MOD;
    }
   
    invFact[MAXN - 1] = modInverse(fact[MAXN - 1]);
    for (int i = MAXN - 2; i >= 0; i--) {
        invFact[i] = (invFact[i + 1] * (i + 1)) % MOD;
    }
}

// 计算 C(n, r)
int nCr(int n, int r) {
    if (r < 0 || r > n) return 0;
    return fact[n] * invFact[r] % MOD * invFact[n - r] % MOD;
}
    \end{lstlisting}
     \subsubsection{线段树}
    \begin{lstlisting}
               struct Segtree {
    // 内部节点结构
    struct Node {
        int l, r;
        long long data;
        long long lazy;
    };

    int n; 
    vector<Node> tr; // 使用 vector 动态开点

    // 局部宏定义，避免污染外部命名空间
    #define ls (u << 1)
    #define rs (u << 1 | 1)

    // 构造函数，初始化大小
    Segtree(int _n) {
        n = _n;
        tr.resize((n + 1) << 2);
    }

    void pushup(int u) {
        tr[u].data = min(tr[ls].data, tr[rs].data);
    }

    void eval(int u, long long v) {
        tr[u].data = min(tr[u].data, v);
        tr[u].lazy = min(tr[u].lazy, v); 
    }

    void pushdown(int u) {
        if (tr[u].lazy != INF) {
            eval(ls, tr[u].lazy);
            eval(rs, tr[u].lazy);
            tr[u].lazy = INF; 
        }
    }

    // 内部递归建树
    void build_tree(int u, int l, int r, const vector<long long>& a) {
        tr[u] = {l, r, INF, INF};
        if (l == r) {
            // 如果传入了有效的初始数组 a，则赋值；否则默认为 INF
            if (l < a.size()) tr[u].data = a[l];
            return;
        }
        int mid = (l + r) >> 1;
        build_tree(ls, l, mid, a);
        build_tree(rs, mid + 1, r, a);
        pushup(u);
    }

    // 内部递归更新
    void update_tree(int u, int L, int R, long long v) {
        if (L <= tr[u].l && R >= tr[u].r) {
            eval(u, v);
            return;
        }
        pushdown(u);
        int mid = (tr[u].l + tr[u].r) >> 1;
        if (L <= mid) update_tree(ls, L, R, v);
        if (R >= mid + 1) update_tree(rs, L, R, v);
        pushup(u);
    }

    // 内部递归查询
    long long query_tree(int u, int L, int R) {
        if (L <= tr[u].l && R >= tr[u].r) {
            return tr[u].data;
        }
        pushdown(u);
        int mid = (tr[u].l + tr[u].r) >> 1;
        long long res = INF; 
        
        if (L <= mid) res = min(res, query_tree(ls, L, R));
        if (R >= mid + 1) res = min(res, query_tree(rs, L, R));
        return res;
    }

    // ==========================================
    // 外部调用的干净接口
    // ==========================================

    // 1. 初始化建树 
    // 如果有初始数组，传进来；如果没有，直接传个空的 vector
    void build(const vector<long long>& a = vector<long long>()) {
        build_tree(1, 1, n, a);
    }

    // 2. 对区间 [L, R] 执行修改 (取 min)
    void update(int L, int R, long long v) {
        if (L > R || L < 1 || R > n) return; // 自动拦截非法区间
        update_tree(1, L, R, v);
    }

    // 3. 查询区间 [L, R] 的结果 (求 min)
    long long query(int L, int R) {
        if (L > R || L < 1 || R > n) return INF;
        return query_tree(1, L, R);
    }

    #undef ls
    #undef rs
};

/*
// 1. 实例化与建树
// 假设你要维护的区间长度是 m
int m = 10000;
SegTree seg(m);

// 情景 A：你需要一个全空的树（比如这道黑题，初始 data 和 lazy 都是 INF）
seg.build(); 

// 情景 B：你有初始数据，需要带数组建树（注意 a 数组有效下标必须从 1 开始匹配）
// vector<long long> arr = {0, 10, 20, 30, 40}; // arr[1] 到 arr[4]
// SegTree seg2(4);
// seg2.build(arr);

// 2. 区间修改与查询
// 区间 [2, 5] 更新值为 100 (由于内部是取 min，只有小于原值才会更新)
seg.update(2, 5, 100);

// 查询区间 [3, 4] 的最小值
long long ans = seg.query(3, 4);

// 查询单点 3 的值（区间左右端点填一样即可）
long long single_val = seg.query(3, 3);
*/
    \end{lstlisting}
    \subsubsection{ST表}
    \begin{lstlisting}
    

/**
 * ST 表模板
 * 复杂度：预处理 O(N log N)，查询 O(1)
 */
const int MAXN = 200005; // 这里的 N 要根据题目总量给够
const int LOGN = 20;     // 2^20 足够覆盖 10^6

int st[MAXN][LOGN];
int lg[MAXN];

// 1. 预处理 log 数组 (在 main 函数开头调用一次即可)
void init_lg() {
    lg[1] = 0;
    for (int i = 2; i < MAXN; i++) {
        lg[i] = lg[i / 2] + 1;
    }
}

// 2. 构建 ST 表 (n 是数组的最大有效下标)
// 注意：如果你的数组是从 0 到 n，那么传参就是 n
void build(int n, const vector<int>& a) {
    // 基础状态：长度为 1 的区间最大值就是本身
    for (int i = 0; i <= n; i++) {
        st[i][0] = a[i];
    }
    
    // 动态规划填充：倍增区间长度
    for (int j = 1; j < LOGN; j++) {
        for (int i = 0; i + (1 << j) - 1 <= n; i++) {
            // 区间 [i, i + 2^j - 1] = max(左半部分, 右半部分)
            st[i][j] = max(st[i][j - 1], st[i + (1 << (j - 1))][j - 1]);
        }
    }
}

// 3. 查询区间 [L, R] 的最大值 (L, R 均为闭区间)
int query(int L, int R) {
    if (L > R) return -1e18; // 异常处理，返回一个极小值
    int k = lg[R - L + 1];
    return max(st[L][k], st[R - (1 << k) + 1][k]);
}
    \end{lstlisting}
   
    \subsubsection{字典树}
    \begin{lstlisting}
    		const int MAXN=1000005;
		int trie[MAXN][26];//trie[u][c]代表节点u经过字符c到达的下一个节点是几号节点
		int cnt[MAXN];//cnt[u]代表有多少个字符串经过了节点u
		int tot=0;//当前节点总数
		void insert(const string& s)//将s插入trie
		{
			int u=0;
			for(char c:s)
			{
				int v=c-'a';
				if(!trie[u][v]) trie[u][v]=++tot;
				u=trie[u][v];
				cnt[u]++;
			}
		}
		ll query(const string& s)//查询s与当前trie中其他字符串的LCP(最长公共前缀长度)之和
		{ 
			int u=0;
			ll ans=0;
			for(char c:s)
			{
				int v=c-'a';
				if(!trie[u][v])
				break;
				u=trie[u][v];
				ans+=cnt[u];
			}
			return ans;
		}
    \end{lstlisting}
    \subsubsection{字符串双哈希}
    \begin{lstlisting}
    		// 核心防御层：基于系统高精度时钟的随机数发生器
mt19937_64 rng(chrono::steady_clock::now().time_since_epoch().count());

// 生成随机进制 (Base)，限制在 [233, 100000] 之间的奇数
long long generate_base() {
    long long b = uniform_int_distribution<long long>(233, 100000)(rng);
    return b % 2 == 0 ? b + 1 : b; 
}

// 全局统一的 Base，确保同一次评测中所有字符串使用相同的计算维度
const long long BASE1 = generate_base();
const long long BASE2 = generate_base();

// 算法核心层：动态随机双哈希结构体
struct DH {
    // 固定的两套冷门大质数
    const long long MOD1 = 1000004323;
    const long long MOD2 = 1000000021;

    int n;
    vector<long long> h1, h2;
    vector<long long> p1, p2;

    // 初始化构造函数
    DH(const string& s) {
        n = s.length();
        // 下标从 1 开始，预留 0 位作为边界，防止越界
        h1.assign(n + 1, 0); h2.assign(n + 1, 0);
        p1.assign(n + 1, 1); p2.assign(n + 1, 1);

        for (int i = 1; i <= n; i++) {
            // 字符映射，加 1 防止 'a'-'a'=0 带来的前导零失效问题
            long long val = s[i - 1] - 'a' + 1; 

            h1[i] = (h1[i - 1] * BASE1 + val) % MOD1;
            h2[i] = (h2[i - 1] * BASE2 + val) % MOD2;

            // 预处理 BASE 的幂次数组，用于 O(1) 提取
            p1[i] = (p1[i - 1] * BASE1) % MOD1;
            p2[i] = (p2[i - 1] * BASE2) % MOD2;
        }
    }

    // O(1) 提取子串 [L, R] 的双哈希值 (1-based)
    pair<long long, long long> get_hash(int L, int R) {
        long long res1 = (h1[R] - h1[L - 1] * p1[R - L + 1]) % MOD1;
        long long res2 = (h2[R] - h2[L - 1] * p2[R - L + 1]) % MOD2;
        
        // 抹平 C++ 负数取模的语言特性坑
        if (res1 < 0) res1 += MOD1;
        if (res2 < 0) res2 += MOD2;
        
        return {res1, res2};
    }
};

/*int main() {
    // 高压防 Hack 测试演示
    string text = "abacabadabacaba";
    DH dh(text);
    
    // 比较两个 "abacaba"
    auto left_part = dh.get_hash(1, 7);
    auto right_part = dh.get_hash(9, 15);
    
    if (left_part == right_part) {
        cout << "匹配成功！本次运行采用的密钥 - BASE1: " << BASE1 << ", BASE2: " << BASE2 << endl;
    }
    
    return 0;
}*/
    \end{lstlisting}
    \subsubsection{马拉车}
    \begin{lstlisting}
    		// 极致封装的 Manacher 模板
struct Manacher {
    string t;       // 预处理后的字符串 (例如: "$#a#b#a#^")
    vector<int> p;  // 回文半径数组

    // 初始化并直接计算出所有的回文半径
    Manacher(const string& s) {
        // 1. 预处理：插入分隔符和边界符
        t = "$#";
        for (char c : s) {
            t += c;
            t += '#';
        }
        t += '^'; // 结尾加上不同的字符，巧妙省去 while 循环里的越界判断

        int n = t.length();
        p.assign(n, 0);
        int C = 0, R = 0; // C: 当前护城河的中心, R: 护城河的最右边界

        // 2. 核心 O(N) 推导
        // 注意遍历范围：跳过开头 '$' 和结尾 '^'
        for (int i = 1; i < n - 1; i++) {
            // 如果 i 在保护伞 R 内部，直接白嫖镜像点 (2*C - i) 的半径
            if (R > i) {
                p[i] = min(R - i, p[2 * C - i]);
            } else {
                p[i] = 0;
            }
            // 尝试暴力向外扩展（如果匹配，半径增大）
            while (t[i + 1 + p[i]] == t[i - 1 - p[i]]) {
                p[i]++;
            }
            // 如果新扩展出的回文串越过了原来的右边界 R，更新护城河
            if (i + p[i] > R) {
                C = i;
                R = i + p[i];
            }
        }
    }

    // 常用接口 1：获取整个字符串的最长回文子串的【长度】
    int get_max_len() {
        int max_len = 0;
        for (int radius : p) {
            max_len = max(max_len, radius);
        }
        return max_len;
    }

    // 常用接口 2：判断原串中闭区间 [L, R] 是否为回文串 (L, R 为 0-based 原串下标)
    bool is_palindrome(int L, int R) {
        // 原串下标映射到处理后的串 t 中的中心位置
        // 原串的第 L 个字符在 t 中的位置是 2*L + 2
        // 原串的第 R 个字符在 t 中的位置是 2*R + 2
        int center = (2 * L + 2 + 2 * R + 2) / 2;
        int len = R - L + 1;
        // 如果该中心的最大回文半径 >= 所需长度，则说明该区间是回文串
        return p[center] >= len;
    }
};

/*int main() {
    string s = "abacabad";
    Manacher man(s);
    
    cout << "原串: " << s << endl;
    cout << "最长回文子串长度: " << man.get_max_len() << endl; // 输出 7 (abacaba)
    
    // O(1) 查询任意区间是否为回文串
    cout << "区间 [0, 2] (aba) 是否回文: " << (man.is_palindrome(0, 2) ? "是" : "否") << endl;
    cout << "区间 [0, 3] (abac) 是否回文: " << (man.is_palindrome(0, 3) ? "是" : "否") << endl;

    return 0;
}*/
    \end{lstlisting}
    \subsubsection{KMP}
    \begin{lstlisting}
    		struct KMP {
    string p;
    vector<int> nxt;

    // 构造函数：传入模式串，自动计算并初始化 next 数组
    KMP(const string& pattern) : p(pattern) {
        int m = p.length();
        nxt.assign(m, 0);
        if (m == 0) return;
        
        for (int i = 1, j = 0; i < m; i++) {
            while (j > 0 && p[i] != p[j]) {
                j = nxt[j - 1];
            }
            if (p[i] == p[j]) {
                j++;
            }
            nxt[i] = j;
        }
    }

    // 匹配函数：传入主串，返回所有匹配成功的起始下标（0-indexed）
    vector<int> match(const string& s) {
        vector<int> res;
        int n = s.length(), m = p.length();
        if (m == 0 || n < m) return res;

        for (int i = 0, j = 0; i < n; i++) {
            while (j > 0 && s[i] != p[j]) {
                j = nxt[j - 1];
            }
            if (s[i] == p[j]) {
                j++;
            }
            if (j == m) {
                res.push_back(i - m + 1);
                j = nxt[j - 1]; // 继续寻找下一个匹配
            }
        }
        return res;
    }
};

/*int main() {
    string text = "ABABABA";
    string pattern = "ABA";

    // 1. 用模式串初始化 KMP 结构体（内部自动计算 next 数组）
    KMP kmp(pattern);

    // 2. 在主串中进行匹配
    vector<int> result = kmp.match(text);

    // 打印结果：0 2 4
    for (int idx : result) {
        cout << idx << " ";
    }
    cout << "\n";
}*/

    \end{lstlisting}
     \subsubsection{关于pq里存node怎么写大小根堆}
    \begin{lstlisting}
    	bool operator<(const node& other)const{
			return d>other.d;
		}
        如果小根堆里面就写大于号，大根堆就写小于号，反着来
    \end{lstlisting}
     \subsubsection{LIS和LNIS}
    \begin{lstlisting}
    		//LNIS
			vector<int>dp(n+1,0);	
			int len1=1;
			dp[1]=a[1];
			for(int i=2;i<=n;i++){
				int pos=upper_bound(dp.begin()+1,dp.begin()+len1+1,a[i],greater<int>())-dp.begin();
				 if(pos==len1+1){
					len1++;
					dp[len1]=a[i];
				 }
				 else{
					dp[pos]=a[i];
				 }
			}
			cout<<len1<<"\n";
			//LIS
			vector<int>dp1(n+1,0);
			dp1[1]=a[1];
			int len2=1;
			for(int i=2;i<=n;i++){
				int pos=lower_bound(dp1.begin()+1,dp1.begin()+len2+1,a[i])-dp1.begin();
				if(pos==len2+1){
					len2++;
					dp1[len2]=a[i];
				}
				else{
					dp1[pos]=a[i];
				}
			}
			cout<<len2<<"\n";
            //这里的dp[i]代表的是长度为i的递增子序列所能达到的最小结尾元素
    Dilworth定理：把一个序列划分成最少数量的“不上升子序列”，其最少数量恰好等于原序列的“最长上升子序列（LIS）”的长度。
    反之同理：把一个序列划分成最少数量的“上升子序列”，等于“最长不上升子序列（LNIS）”的长度。

    对于递减数组，upper_bound是找第一个小于的，lower_bound是找第一个小于等于的，然后注意递减数组写的时候要最后加一个greater<int>().
    \end{lstlisting}
    \subsubsection{bitset}
    \begin{lstlisting}
    bitset是一个c++标准库的特殊容器，可以看作是一个固定长度的、超级巨大的二进制整数。
    用法：在全局写bitset<长度>a，长度到一亿都没问题，而且注意必须得在全局写。
    a[0]=1,含义就是把第0位改成1，第0位是最右边的个位。
    如果要开多个bitset，就是开个bitset二维数组的话，就是bitset<长度>a[数组个数]。
    \end{lstlisting}
    \section{数据结构}
    \subsection{trick}
    \subsubsection{关于树状数组的感悟}
    \begin{lstlisting}
        树状数组最常见的应用场景就是：我正在从左往右(从右往左)扫描，需要快速统计之前出现过的某类数有多少个。

        树状数组的下标，通常就是你想要比较大小的那个属性。比如经典的求逆序对，我遍历到i时我想知道前面有多少个数大于a[i],那这时候我们想要比较大小的属性显然就是a[i]对吧，所以把a[i]作为树状数组的下标，每次add(a[i],1);这里后面的加什么当然是根据题目来看了。

        面对一道树状数组题时，先问自己：当我遍历到当前位置时，我想知道前面的哪些点，与当前点满足什么关系？这些点的某种属性的统计值是多少？把这句话理清楚了就知道怎么写add了。

        如果要建多个树状数组，比如优化二维dp的时候，把tree改成vector<vector<int>>就行了，然后下面的add，query啥的多个参数就行，表示查询第几个树状数组。

        树状数组和dp其实关系非常密切。因为你做dp的时候也经常会出现：当前状态=若干历史状态的贡献之和。而树状数组干的事情就是动态维护一批数的前缀和并快速查询。所以用树状数组优化dp是很常见的技巧。
    \end{lstlisting}
    \subsubsection{一个树状数组经典小trick}
    \begin{lstlisting}
        有的题需要我们“动态维护元素相对顺序”，比如把一个数放到最前面，在他前面的数依次向后移一位这种操作，如果真的要模拟那就是O(n)，多次操作就爆炸了。那么怎么办呢？
        答案是：假设有n个数，m个操作，就开一个长度为n+m的数组。先把所有数放在数组的最右边，然后每次对一个数操作时，就直接把这个数放到当前第一个数的左边，然后他原来的位置就直接空着就好了。因为每次操作最多会向左拓展一个位置(除了对第一个数操作不会向左拓展)，所以m次操作最多拓展m个位置，这就是为什么把数组开成n+m的原因。
        这样操作的话，我们每次只需要动一个数就好了，不需要牵一发而动全身的去动另外的数，然后后面要进行什么查询可以用树状数组完成，比如查询每个元素在移动过程中到达过的最大最小位置什么的。
    \end{lstlisting}
    \subsection{板子}
    \subsubsection{并查集} % 小标题
    \begin{lstlisting}
        //并查集的所有操作都是近似O(1)的，不用怕t
        //并查集可以维护每个连通块的最大最小值，加个mx,mn数组，先初始化为每个点自己的值，merge的时候更新一下就行了。
        struct dsu {
    vector<int> fa, sz;
    int cnt; // 当前集合数量（连通块数量）

    dsu(int n = 0) {
        init(n);
    }

    // 初始化
    void init(int n) {
        fa.resize(n + 1);
        sz.assign(n + 1, 1);
        iota(fa.begin(), fa.end(), 0);
        cnt = n;
    }

    // 查找根节点 + 路径压缩
    int find(int x) {
        if (fa[x] == x) return x;
        return fa[x] = find(fa[x]);
    }

    // 合并两个集合
    // 返回 true 表示成功合并
    // 返回 false 表示原本就在同一个集合
    bool merge(int x, int y) {
        x = find(x);
        y = find(y);

        if (x == y) return false;

        // 按大小合并
        if (sz[x] < sz[y]) swap(x, y);

        fa[y] = x;
        sz[x] += sz[y];

        cnt--;
        return true;
    }
    // 判断是否属于同一个集合
    bool same(int x, int y) {
        return find(x) == find(y);
    }
    // 查询 x 所在集合的大小
    int size(int x) {
        return sz[find(x)];
    }
};
    \end{lstlisting}
    \subsubsection{树状数组}
    \begin{lstlisting}
        //注意树状数组必须是1-base，就是你不能add(0,1)，这样会死循环。
        //注意如果数据范围达到1e9这种，树状数组开不了这么大，可以离散化。
        //注意如果用树状数组处理前缀和的话，不要忘记考虑pre[i]-pre[0]的情况了，就是说别忘记把pre[0]也add到树状数组里面。
        struct BIT {
    int n;
    vector<long long> tree;

    BIT(int n) {
        init(n);
    }

    void init(int n) {
        this->n = n;
        tree.assign(n+1,0);
    }

    int lowbit(int x) {
        return x & -x;
    }

    // 单点增加
    void add(int x,long long k) {
        while(x<=n) {
            tree[x]+=k;
            x+=lowbit(x);
        }
    }

    // 查询 [1,x]
    long long query(int x) {
        long long ans=0;

        while(x>0) {
            ans+=tree[x];
            x-=lowbit(x);
        }

        return ans;
    }

    // 查询 [l,r]
    long long Query(int l,int r) {
        return query(r)-query(l-1);
    }
};
    \end{lstlisting}
    \section{图论}
    \subsection{trick}
    \subsubsection{MST与BST的关系，瓶颈问题}
    \begin{lstlisting}
    MST：最小生成树  BST：瓶颈生成树
    瓶颈生成树的定义：所有生成树中，最大边权最小的那棵树
    **任何图的MST，必定是BST。//但是BST不一定是MST
    只要看到题目有“\textbf{使得所有路径的最大边权最小}”“\textbf{...最大距离最小}”直接写MST
    对于图上任意两点之间“最大边权最小”的路径是mst上的路径。
    \end{lstlisting}
    \subsubsection{最大生成树}
    \begin{lstlisting}
    要得到一个图的最大生成树，其实最简单的方法就是对图的每个边权先取反，然后求最小生成树就行了。
    与mst“最大边权最小”的性质一样，最大生成树有“最小边权最大”的性质。
    \end{lstlisting}
    \subsubsection{聚类问题}
    \begin{lstlisting}
    聚类问题就是说：有一堆离散的点，你要把他们分成k个连通块，就是聚在一起嘛，然后要满足一些条件，比如连通块之间最大距离最小，或者连通块内部的最小距离最大这样子。
    这类问题其实挺典的吧我觉得，代码都不用怎么改。
    我现在见过的就四种情况，其实就是连通块内部还是外部，最大距离最小还是最小距离最大排列组合得到的。
    
    如果是块外部，那么相当于刚开始的“候选人”是所有边，然后你每连一条边就相当于把它从候选人里去掉。如果你是要求最大距离最小，那想一想，一个升序数组，你要去掉一些值，使剩下的最大值最小，那肯定是从大到小去对吧，那搞个最大生成树把边排下序取其中第n-k项什么的（具体问题具体分析）就行了。然后如果是最小距离最大，那就是最小生成树。

    如果是块内部，那么相当于刚开始得候选人是空集，你每连一条边就是把这条边加到候选人里面。如果你要求最大距离最小，那肯定是从小往大加边对吧，那就是搞最小生成树。反过来一样的，用最大生成树即可。

    聚类问题一般都是给坐标，就是隐式图形式，所以用prim解决（至少我现在只见过坐标形式，可能还有其他的但是我没见过）。然后用prim的过程中有些注意事项，详情见板子的prim部分。
    \end{lstlisting}
    \subsubsection{拓扑排序接dp（DAG上的dp）}
    \begin{lstlisting}
    对一个DAG拓扑排序后，对得到的线性序列做dp是很常用的技巧，因为拓扑排序完美解决了dp最苛刻的前提条件——无后效性。
    你站在节点u准备更新别人时，所有能影响u的前置节点一定都已经计算完毕了。
    而且DAG上dp往往用“刷表法”，意思就是站在u去更新后面的在adj[u]中的v，与之相对的“填表法”意思就是准备更新v，发现前面的u还没算，于是先算u再把u的值拉回来。在dag上做dfs用填表法会更自然，做dp用刷表法会更顺畅。
    dagdp可以解决很多问题，例如dag上的最短路最长路，从起点到终点有多少种路径，图上博弈论，附带背包或多维度的条件极值...
    \end{lstlisting}
    \subsubsection{scc->dagdp}
    \begin{lstlisting}
    因为scc缩点之后得到的就是一个dag，所以scc接dagdp是非常常见且固定的套路。具体操作就是scc完之后遍历每个点的每条边，如果u，v的sccid不相等就连条边在sccid[u],sccid[v].dag就建好了，然后后面就正常dagdp。
     vector<pair<int, int>> edges; // 拿个小本本把边记下来
            
            for(int i=1; i<=n; i++){
                for(auto x: g.adj[i]){
                    int u = g.scc_id[i];
                    int v = g.scc_id[x];
                    if(u != v){
                        edges.push_back({u, v}); // 先不急着建图，全存起来
                    }
                }
            }
            
            // 核心魔法：排序 + 去重
            sort(edges.begin(), edges.end());
            edges.erase(unique(edges.begin(), edges.end()), edges.end());
            topo ts(g.scc_cnt);
            // 现在 edges 里面绝对没有重边了，安心建图
            for(auto e : edges){
                ts.add_edge(e.first, e.second);
            }
            //这就是scc之后建dag的代码，因为可能有重边所以需要去重
    \end{lstlisting}
    \subsubsection{dag至少加几条边能变成一个scc?}
    \begin{lstlisting}
    一个dag至少加max(入度为0的点的个数，出度为0的点的个数)能变成一个scc，注意一个特判是如果只有一个点那答案是0。
    只要你入度为0，那么想要变成scc，是必须从其他点连一条边到这个点的，同样的出度为0也必须从这个点连个边到其他点才有可能变成scc。相当于加边个数是>=max(入度0个数，出度0个数)，相当于证明了必要性吧。然后，你只要把入度为0和出度为0的两点之间连边，就一定能在只给入度/出度为0的点连边的情况下使其变成scc，相当于证明了充分性 
    但是你想要构造的话有点麻烦，主要想法是错位着去连，就是如果两个点之间已经有一条边，那尽量避免去再连一条，因为这样只能形成一个两个点的小环，不优。
    \end{lstlisting}
    \subsubsection{题目没说有多少个点并且点不一定连续怎么建图初始化}
    \begin{lstlisting}
        如果告诉你一个节点个数上限，比如1000，2e5这样比较小的，可以先初始化为g(1005),g(200005)这样，然后在struct里建一个exist布尔数组，在add_edge的时候把出现的边在exist里标记为true，跑tarjan的时候只遍历存在的点。
    \end{lstlisting}
    \subsubsection{记录scc内部边权之和}
    \begin{lstlisting}
        其实看到这个问题我第一眼想法是在tarjan内部dfs过程中怎么记录一下边权，但其实没必要，直接跑完scc之后遍历每条边，然后如果这条边的两个端点属于同一个scc，就把对应的scc权值和加上这条边的权值。事先创个vector<node>edge，读入每条边的时候记录一下就行。
    \end{lstlisting}
    \subsubsection{求dag上特定起点最长路}
    \begin{lstlisting}
        求dag上最长路显然用dagdp，但是如果有特定起点的话呢，这里要注意拓扑序不能表示连通关系，所以直接从起点处开始dp是不行的，因为后面的点可能根本就跟起点不连通。所以采用的方法是把dp值都初始化为-INF，只把起点赋值。遍历到一个点如果他的dp值是-INF那么代表起点到不了他，直接continue。这样所有dp值被更新的点都是和起点连通的点，不连通的一直保持-INF。其实就相当于搞了个vector<bool>vis标记一下，true代表可达，false代表不可达。
    \end{lstlisting}
    \subsubsection{判断无向图中有没有环,以及看这个环是几元环}
    \begin{lstlisting}
        看无向图有没有环直接从任意起点开始dfs遍历一遍整个图就行了，如果遍历到一个节点，去看他连的其他点，如果有一个点已经被访问过且不是他的父节点，那肯定就是出现了环。
        然后看是几元环的话也很简单，就是加个深度数组vector<int>dep，在dfs过程中顺手给每个点标记一下深度，然后在发现一个已经访问过的节点v时，dep[u]-dep[v]+1就是边数。
    \end{lstlisting}
    \subsubsection{一些关于scc的感悟}
    \begin{lstlisting}
        看到形如“只要u怎么样，v就必须怎么样”这种约束关系，就可以尝试u->v这样连边然后看怎么用scc可以搞一下，因为这些约束关系就导致了他们是相互依赖无法拆分的整体，自然就想到缩点。
        看到形如“我要从一堆点的集合中选出一个子集，要求子集中的每个点的每条边都必须连向子集内的点，即不能有出边连向不在子集内的点，再概括的说就是子集对于出边封闭”，那么这时候我们也可以很自然的想到scc。为什么呢，因为我们想象对一个图scc之后建的dag，这时候每个scc内部就是出边封闭的。它的出边都体现在dag的连边上。所以如果我们想要找到一个出边封闭的集合，那出度=0的scc里的所有点就是一个合法集合。
        然后就是，scc不仅可以作为一种分析工具，就是我有一些性质，然后我自然想到用scc去缩点，简化问题，还可以作为一种构造工具，就是说要求我构造出一个合法的集合，然后我尝试scc一下看看会发生什么，然后发现，欸，如果我直接把出度为0的scc里的所有点作为我选的集合，不就是一个成功的构造吗？

        有向图中只要两个点相互可达，那么他们一定在一个scc中；只要两个点不相互可达，那么他们一定不在一个scc中。这个问题看起来很蠢，但是我在学2-sat的时候真的问了这样的问题；“有没有可能两个点相互可达但是不在同一个scc中？”可以用反证法来想，假设有两个点真的相互可达但是他们不在一个scc中，那么他们俩所在的scc之间就会连一条双向边，但scc缩完点必然是一个dag，矛盾。或者直接从定义去想，如果两个点相互可达，那么不管其他点，他们俩至少能形成一个两个点的scc。
        这个结论换个话说就是，如果要判断两个点是否相互可达，那么看他们是不是在同一个scc中就行了，这是充要条件。这也就是2-sat的核心结论。

        缩点dag中如果有边a->b,那么scc_id[a]>scc_id[b]，也就是说tarjan分配scc_id的顺序是逆拓扑序。注意这里是“每个scc弹栈时编号递增写法”就会这样，换个写法可能就相反了。
        也就是说，如果缩点dag中，a~>b(意思是a可达b)，那么scc_id[a]>scc_id[b],但是注意反过来并不成立，因为scc_id是一定要有大小关系的，就算你的a和b并没有直接相连关系(也就是并列关系)，也不能两个id相等，一定会有一个大小关系。
    \end{lstlisting}
    \subsubsection{有环的连通块删一条边后不一定依旧保持连通}
    \begin{lstlisting}
        这个问题看起来很蠢很显然，但是我做题的时候就想当然了。
        有环的连通块的边，可以分为在环上的边，和不在环上的边，也就是桥。你删掉在环上的边会依旧保持连通，但是删掉桥会割成两部分。
        那要怎么知道，或者说在连边的时候记录在环上的边呢，其实很简单，就是用并查集，u和v连边时merge(u,v),如果发现u和v已经在同一个连通块了那肯定这条边就是在环上的边对吧。
    \end{lstlisting}
    \subsubsection{拓扑排序如果有自环也会触发判断有环}
    \begin{lstlisting}
        没什么好说的，注意一下就行了。
    \end{lstlisting}
    \subsubsection{dagdp可以跟bitset搭配使用}
    \begin{lstlisting}
        dagdp有时候是可以跟bitset一起搭配使用的，比如要你求一个有向图中每个点可以到达的点数，那就可以搞一个bitset数组，对每个点都搞一个bitset，用来记录他可以到达哪些点，然后dp转移就是dp[u]|=dp[v]，意思就是有条边u->v，那么v能到的点u都能到，所以|一下，这也是bitset的合并的通用方式吧。
        还有一个点就是如果点太多，或者说点多才是常见情况吧，比如n=2e5，那你要开个2e5*2e5的二维数组(二维bitset数组的内存占用可以看作二维数组，但是一维bitset可以开非常大)，那就爆爆的爆了。所以我们需要分块bitset。
        分块bitset的意思就是，每次只处理少量点，处理多次把答案加起来。比如n=5e4的话，你每次处理64个数，处理700多次就行了，时间上64*700非常小，空间上2e5*64也非常小。比如上面那个例题，就是你每次去看所有点能否到达你给定的64个点，然后处理700多次把每次能到达的点数加起来。
        分块bitset的精髓在哪呢，就是我们的问题就是你编号太大的话装不下对吧，比如来个10000编号，你没法开那么大的bitset数组，那么你每次只处理64个点，你把他们的编号都映射到[0,63]上面，比如你要处理[256,319],就映射成[0,63]。每次只用开64长度的bitset数组。
    \end{lstlisting}
    \subsubsection{对于无向连通图，删k条边最多增加k个连通块}
    \begin{lstlisting}
        这个结论其实很显然吧，你删一条边最多只能把一个连通块分成两个，也就是增加一个，那删k条边就最多增加k个。
        而且你试想一下，怎么可能删一条边把一个连通块分成三个呢？
    \end{lstlisting}
    \subsubsection{怎么求一个无向图删掉一条桥后，分成的两个连通块各自的节点数？}
    \begin{lstlisting}
        其实这个也很简单，只要在tarjan里加一个数组sz,sz[i]表示i节点的子树大小，然后你割掉一个桥，儿子一侧的子树大小就是其中一个块的数量对吧，然后另一块当然就是用n-这一块(注意这个是原图是连通的情况，如果不连通那应该用这个连通块的大小-这一块)。
        然后注意必须是儿子一侧的大小，如果是父亲一侧的子树大小那是没意义的，不是任何一块的大小，自己想象一下就知道了。所以我们在遍历到一个桥的时候，按照我们之前的想法就是把sz[儿子]取出来，那么问题来了，怎么知道哪边是儿子呢？显然如果要等到调用的时候再去找谁是儿子是很困难的(后来才发现其实并非困难，所以下面的方法也不是很必要了。。详情见下一节。。)，所以不妨再设一个bsz(bridge_sz)数组，记录每个桥如果被割能形成的其中一个连通块的大小。具体操作就是：直接在dfs发现一条桥的时候就把这个桥的bsz[id]设为sz[v]，这个v是儿子。
        然后我们遍历到一个桥直接取这个桥的bsz就行了。记录bsz的代码如下：
                if (low[v] > dfn[u]) {
                    br[id] = 1;
                    bsz[id]=sz[v];
                }
        哦对了，还没说怎么求sz，这个也很简单，就是先把所有点的sz都初始化为1，因为他们初始只有自己一个点嘛，然后dfs过程中对于一个u对他的所有儿子v进行sz[u]+=sz[v]就行了，很容易理解。
        能这样做的一个关键点在于，你在发现一个桥的时候，是已经把儿子节点v，或者说儿子这一侧的子树的所有节点的low，sz什么的都算好了的，所以才能直接bsz[id]=sz[v]，直接记录这个桥的bsz，我是这么想的。
        
        然后还有一个要注意的点是，或者说我们需要理解的点是，\textbf{桥两侧大小是图本身的性质，而不是dfs树的性质，dfs只是帮助我们计算的工具}。什么意思呢，就是你不管以哪个点为起点进行dfs，不论你形成的dfs树是什么形态，图本身的性质并不会因此而改变，也就是说桥并不会改变，原来是桥，并不会因为你换个dfs顺序就变成不是桥，因为图本身没变，而且你割掉一条桥后两边的块大小也不会改变的，这个非常好理解对吧。而dfs树只是我们求解它的方式，我们以不同的起点来dfs，构建出来的dfs树长得不一样，但是求出来的答案肯定都是对的。影响的只是，你在一个桥上哪边是儿子节点，这个桥的bsz取哪边节点的sz。(这么想其实很容易发现一个桥的bsz在不同的dfs树下最多只能有两个取值，对应两边的子树大小。没什么用但是有助于增进理解吧我个人感觉。)
    \end{lstlisting}
    \subsubsection{怎么快速知道一个无向图中的一条边，在dfs树中谁是儿子？}
    \begin{lstlisting}
        如果是比较一条树上边，那么直接比较他们的dfn就行了，dfn大的肯定就是儿子。
        如果是一条非树边，那么就不能直接比较dfn了，要更复杂一点，具体我还没弄清楚。
    \end{lstlisting}
    \subsubsection{无向图中桥必定在dfs树上}
    \begin{lstlisting}
        这个也很显然。因为你桥割了那整个图都不连通了，那dfs树怎么可能不经过桥就能遍历完整张图呢对吧。
    \end{lstlisting}
    \subsubsection{对于一个无向图，在环上的边一定不是桥,不在环上的边一定是桥}
    \begin{lstlisting}
        这个结论也很显然，前面其实也提到过，但还是单独拿出来说一下吧。
        你删掉一个在环上的边，肯定不会破坏原图的连通性；反之，删掉一条不在环上的边，肯定会破坏连通性。很好理解。
    \end{lstlisting}
    \subsubsection{做无向图的tarjan的时候，发现一个dfn[v]不等于0，就是发现一个点已经去过了，不能确定v是祖先还是u是祖先}
    \begin{lstlisting}
        这个我自己反正是之前没想清楚的，我以为你发现一个dfn[v]不是0，那肯定代表发现一条返祖边了，所以肯定v是祖先。但其实不然。
        想象一个最简单的三元环就行了，1连2，2连3，3连1，然后我们dfs树是123，3到1是一个返祖边。那么你站在3去看1的时候，确实是条返祖边，但是你要意识到，1是通过2再到3的，也就是说3弄完了其实还有1->3这条边是没走的，还要回溯回去站在1看1-3这条边，那这时候就发现了dfn[3]不等于0啊，但实际上这不是一条返祖边，它是一条前向边。

        说到这前向边的定义其实就是从祖先到儿子的非树边，它跟返祖边是对应的，只是方向不同。(其实无向图中一般不称前向边，只有树边和返祖边，也没有横叉边，然后有向图有四种边，有树边返祖边前向边横叉边，我只是为了定义无向边的方向所以称其前向边返祖边)
    \end{lstlisting}
    \subsection{板子}
    \subsubsection{拓扑排序}
    \begin{lstlisting}
        struct TopoSort {
    int n;
    vector<vector<int>> adj;
    vector<int> in_deg;
    vector<int> order; // 存储最终的拓扑序列

    // 构造函数，_n 为节点总数。默认开 _n+1 的空间，兼容编号 1~n 或 0~n-1
    TopoSort(int _n) {
        init(_n);
    }

    // 初始化/重置，适合多组测试数据
    void init(int _n) {
        n = _n;
        adj.assign(n + 1, vector<int>());
        in_deg.assign(n + 1, 0);
        order.clear();
    }

    // 添加一条从 u 到 v 的有向边
    void add_edge(int u, int v) {
        adj[u].push_back(v);
        in_deg[v]++;
    }

    // 执行拓扑排序。返回 true 表示图无环且排序成功；返回 false 表示图中存在环
    bool build() {
        queue<int> q;
        // 注意：如果你是 0-indexed，这里的循环应该是 i = 0; i < n; i++
        // 这里默认采用最常见的 1-indexed (1 到 n)
        for (int i = 1; i <= n; i++) {
            if (in_deg[i] == 0) {
                q.push(i);
            }
        }

        order.clear();
        while (!q.empty()) {
            int u = q.front();
            q.pop();
            order.push_back(u);

            for (int v : adj[u]) {
                if (--in_deg[v] == 0) { // 入度减 1，若减为 0 则入队
                    q.push(v);
                }
            }
        }
        
        // 如果序列中的节点数等于图的总节点数，说明无环
        return order.size() == n; 
    }
};

/*void solve() {
    int n, m; 
    cin >> n >> m; // n 个点，m 条边

    // 1. 实例化拓扑排序对象
    TopoSort ts(n);

    // 2. 建图
    for (int i = 0; i < m; i++) {
        int u, v;
        cin >> u >> v;
        ts.add_edge(u, v);
    }

    // 3. 执行算法并判断是否有环
    if (ts.build()) {
        cout << "拓扑序为: ";
        for (int u : ts.order) {
            cout << u << " ";
        }
        cout << "\n";
    } else {
        cout << "图中存在环，无法进行拓扑排序！\n";
    }
}
*/
    \end{lstlisting}
    \subsubsection{dagTR(传递约简)}
    \begin{lstlisting}
        //注意这个传递约简复杂度有点复杂，简要记一下是5000通常可以，10000就有点极限了，可以尝试一下，再往上就最好别考虑了。
struct dagTR{
    using ull = unsigned long long;

    int n;
    topo tp;
    vector<vector<int>> res;

    dagTR(int _n) : n(_n), tp(_n) {}

    // 添加有向边
    void add_edge(int u, int v) {
        tp.add_edge(u, v);
    }

    // 执行传递约简
    // 返回 false 表示原图不是 DAG
    bool build() {
        res.assign(n + 1, {});

        // 1. 拓扑排序
        if (!tp.build()) return false;

        vector<int> pos(n + 1);
        for (int i = 0; i < n; i++) {
            pos[tp.order[i]] = i;
        }

        // 2. 求每个点的可达集合
        int W = (n + 64) / 64;

        vector<vector<ull>> reach(
            n + 1, vector<ull>(W, 0)
        );

        for (int i = n - 1; i >= 0; i--) {
            int u = tp.order[i];

            reach[u][u >> 6] |= 1ULL << (u & 63);

            for (int v : tp.adj[u]) {
                for (int j = 0; j < W; j++) {
                    reach[u][j] |= reach[v][j];
                }
            }
        }

        // 3. 传递约简
        for (int u = 1; u <= n; u++) {
            auto adj = tp.adj[u];

            sort(adj.begin(), adj.end(),
                [&](int a, int b) {
                    return pos[a] < pos[b];
                }
            );

            vector<ull> vis(W, 0);

            for (int v : adj) {
                // 已经能够到达 v，说明这条边冗余
                if ((vis[v >> 6] >> (v & 63)) & 1ULL) {
                    continue;
                }

                // 保留这条边
                res[u].push_back(v);

                // 更新可达集合
                for (int j = 0; j < W; j++) {
                    vis[j] |= reach[v][j];
                }
            }
        }

        return true;
    }
};
/*用法示例
//注意要把struct topo放在这个dagTR上面，这个板子是依赖于topo的。
int n, m;
cin >> n >> m;
dagTR tr(n);

for (int i = 1; i <= m; i++) {
    int u, v;
    cin >> u >> v;
    tr.add_edge(u, v);
}

if (!tr.build()) return;

auto &g = tr.res;
// g 就是传递约简后的邻接表

// 正常遍历邻接表 
for (int u = 1; u <= n; u++) {
    for (int v : g[u]) {
        cout << u << " -> " << v << '\n';
    }
}
*/
    \end{lstlisting}
     \subsubsection{欧拉回路}
    \begin{lstlisting}
    void hierholzer(int u)
{
    stack<int>sk;
    sk.push(u);
    vector<int>path;
    while(!sk.empty())
    {
        int t=sk.top();
        if(adj[t].empty())
        {
            path.pb(t);
            sk.pop_back();
        }
        else
        {
            int v=adj[t].back();
            adj[t].pop_back();
            sk.push(v);
        }
    }
}
    \end{lstlisting}
      \subsubsection{Dijkstra}
    \begin{lstlisting}
    		#include<bits/stdc++.h>
		using namespace std;
		using ll = long long;
		#define pb push_back
		#define pii pair<int,int>
		#define int long long
		#define i128 __int128
        #define ull unsigned long long
        const int INF=2e18;

        
    struct Dij {
    

    struct Edge {
        int to;
        long long weight;
    };

    int n; // 节点数
    vector<vector<Edge>> adj; // 邻接表
    vector<long long> dist;   // 距离数组
	vector<int>pre;
	vector<int>path;
    // 构造函数，初始化大小
    Dij(int _n) {
        n = _n;
        adj.resize(n + 1);
        dist.assign(n + 1, INF);
		pre.assign(n+1,0);
    }

    // 添加无向边
    void add_edge(int u, int v, long long w) {
        adj[u].push_back({v, w});
        adj[v].push_back({u, w}); 
    }

    // 运行 Dijkstra 算法求解单源最短路
    void run(int start) {
        // 如果对同一个图多次查询，需要每次重置距离数组
        dist.assign(n + 1, INF); 
        dist[start] = 0;
		pre.assign(n+1,0);
        priority_queue<pair<long long, int>, vector<pair<long long, int>>, greater<pair<long long, int>>> pq;
        pq.push({0, start});

        while (!pq.empty()) {
            auto [d, u] = pq.top();
            pq.pop();

            if (d > dist[u]) continue;

            for (const auto& edge : adj[u]) {
                int v = edge.to;
                long long weight = edge.weight;

                if (dist[u] + weight < dist[v]) {
                    dist[v] = dist[u] + weight;
                    pq.push({dist[v], v});
					pre[v]=u;
                }
            }
        }
    }

    // 获取起点到目标点的最短距离
    long long get_dist(int target) {
        return dist[target];
    }
	void get_path(int n){
		int cur=n;
		while(pre[cur]!=0){
			path.pb(cur);
			cur=pre[cur];
		}
		path.pb(1);
		reverse(path.begin(),path.end());
	}
};

/*int main() {
    // 实例化结构体，传入节点数量
    Dijkstra graph(n);
    // 读入边
    for (int i = 0; i < m; ++i) {
        int u, v;
        long long w;
        cin >> u >> v >> w;
        graph.add_edge(u, v, w);
    }
    int start_node = 1;
    // 运行算法
    graph.run(start_node);

    // 获取结果
    long long ans = graph.get_dist(n);
    if (ans == graph.INF) {
        cout << -1 << "\n";
    } else {
        cout << ans << "\n";
    }
}*/
    \end{lstlisting}
    \subsubsection{Johnson}
    \begin{lstlisting}
    		struct Johnson {
    const long long INF = 1e18; // 极大值

    struct Edge {
        int to;
        long long weight;
    };

    int n; 
    vector<vector<Edge>> adj; // 邻接表
    vector<long long> h;      // 势能数组
    vector<long long> dist;   // 单次查询的真实距离数组

    // 构造函数，初始化大小
    Johnson(int _n) {
        n = _n;
        // 分配 n+1 的空间，0号点作为超级源点保留
        adj.resize(n + 1);
        h.assign(n + 1, INF);
        dist.assign(n + 1, INF);
    }

    // 添加有向边（带负权的图通常是有向图，若存在无向负权边则必然是负权环）
    void add_edge(int u, int v, long long w) {
        adj[u].push_back({v, w});
    }

    // 步骤 1：运行 SPFA 计算势能，返回 false 表示存在负权环
    bool build_potential() {
        h.assign(n + 1, INF);
        vector<int> cnt(n + 1, 0);
        vector<bool> inq(n + 1, false);
        queue<int> q;

        // 超级源点 0 初始化
        h[0] = 0;
        q.push(0);
        inq[0] = true;

        // 建立超级源点到所有点的 0 权边
        for (int i = 1; i <= n; ++i) {
            adj[0].push_back({i, 0});
        }

        while (!q.empty()) {
            int u = q.front();
            q.pop();
            inq[u] = false;

            for (const auto& edge : adj[u]) {
                int v = edge.to;
                long long weight = edge.weight;

                if (h[u] + weight < h[v]) {
                    h[v] = h[u] + weight;
                    if (!inq[v]) {
                        q.push(v);
                        inq[v] = true;
                        cnt[v]++;
                        // 包含超级源点共 n+1 个节点，松弛次数超标即为负权环
                        if (cnt[v] >= n + 1) return false; 
                    }
                }
            }
        }
        return true;
    }

    // 步骤 2：以 start 为起点运行 Dijkstra
    void run_dijkstra(int start) {
        dist.assign(n + 1, INF);
        dist[start] = 0;

        priority_queue<pair<long long, int>, vector<pair<long long, int>>, greater<pair<long long, int>>> pq;
        pq.push({0, start});

        while (!pq.empty()) {
            auto [d, u] = pq.top();
            pq.pop();

            if (d > dist[u]) continue;

            for (const auto& edge : adj[u]) {
                int v = edge.to;
                long long weight = edge.weight;
                
                // 核心：使用势能修改后的非负边权计算
                long long new_weight = weight + h[u] - h[v];

                if (dist[u] + new_weight < dist[v]) {
                    dist[v] = dist[u] + new_weight;
                    pq.push({dist[v], v});
                }
            }
        }

        // 算法结束后，将计算出的距离还原为真实距离
        for (int i = 1; i <= n; ++i) {
            if (dist[i] != INF) {
                dist[i] = dist[i] - h[start] + h[i];
            }
        }
    }

    // 获取起点到目标点的真实最短距离
    long long get_dist(int target) {
        return dist[target];
    }
};

/*用法示例：
int main() {
    int n = 4, m = 5;
    Johnson graph(n);

    // 1. 建图
    graph.add_edge(1, 2, 2);
    graph.add_edge(2, 3, -4); // 负权边
    // ...读入其余边

    // 2. 检查负权环并构建势能
    if (!graph.build_potential()) {
        cout << "图里存在负权环，无法求解最短路！\n";
        return 0;
    }

    // 3. 对需要的起点跑 Dijkstra
    int start_node = 1;
    graph.run_dijkstra(start_node);

    // 4. 获取结果
    long long ans = graph.get_dist(3); // 获取 1 到 3 的真实距离
    if (ans == graph.INF) {
        cout << "不可达\n";
    } else {
        cout << ans << "\n";
    }

    return 0;
}
*/
    \end{lstlisting}
    \subsubsection{SPFA}
    \begin{lstlisting}
    	struct SPFA {
    const long long INF = 1e18; // 极大值

    struct Edge {
        int to;
        long long weight;
    };

    int n; 
    vector<vector<Edge>> adj; // 邻接表
    vector<long long> dist;   // 距离数组

    // 构造函数，初始化大小
    SPFA(int _n) {
        n = _n;
        adj.resize(n + 1);
        dist.assign(n + 1, INF);
    }

    // 添加有向边
    void add_edge(int u, int v, long long w) {
        adj[u].push_back({v, w});
    }

    // 运行 SPFA，返回 false 表示存在负权环，返回 true 表示正常求出最短路
    bool run(int start) {
        dist.assign(n + 1, INF); 
        vector<int> cnt(n + 1, 0);       // 记录节点入队次数
        vector<bool> inq(n + 1, false);  // 记录节点是否在队列中
        queue<int> q;

        dist[start] = 0;
        q.push(start);
        inq[start] = true;
        cnt[start] = 1; // 起点第一次入队

        while (!q.empty()) {
            int u = q.front();
            q.pop();
            inq[u] = false;

            for (const auto& edge : adj[u]) {
                int v = edge.to;
                long long weight = edge.weight;

                if (dist[u] + weight < dist[v]) {
                    dist[v] = dist[u] + weight;
                    if (!inq[v]) {
                        q.push(v);
                        inq[v] = true;
                        cnt[v]++;
                        // 如果某个点入队次数达到 n 次，说明必然绕了负权环
                        if (cnt[v] >= n) return false; 
                    }
                }
            }
        }
        return true; 
    }

    long long get_dist(int target) {
        return dist[target];
    }
};

/*
int main() {
    int n = 4;
    SPFA graph(n);

    // 建图
    graph.add_edge(1, 2, 2);

    int start_node = 1;
    
    // 运行算法并拦截负权环
    if (!graph.run(start_node)) {
        cout << "图中存在负权环，最短路无解！\n";
    } else {
        long long ans = graph.get_dist(3);
        if (ans == graph.INF) {
            cout << "不可达\n";
        } else {
            cout << "起点到 3 号点的最短距离为: " << ans << "\n";
        }
    }
    return 0;
}*/
	
    \end{lstlisting}
    \subsubsection{Bellman}
    \begin{lstlisting}
    struct BellmanFord {
    const long long INF = 1e18;

    // 直接存储每一条完整的边
    struct Edge {
        int u, v;
        long long weight;
    };

    int n;
    vector<Edge> edges;     // 边集数组
    vector<long long> dist; // 距离数组

    BellmanFord(int _n) {
        n = _n;
        dist.assign(n + 1, INF);
    }

    // 添加有向边
    void add_edge(int u, int v, long long w) {
        edges.push_back({u, v, w});
    }

    // 运行算法，k 表示“最多允许经过 k 条边”
    void run(int start, int k) {
        dist.assign(n + 1, INF);
        dist[start] = 0;

        // 严格限制外层循环跑 k 轮
        for (int i = 0; i < k; ++i) {
            // 核心防御：备份上一轮的距离，防止本轮发生越级“串联”更新
            vector<long long> backup = dist; 
            bool relaxed = false; // 剪枝优化标记

            for (const auto& edge : edges) {
                int u = edge.u, v = edge.v;
                long long w = edge.weight;

                // 必须使用 backup[u] 去更新别人
                if (backup[u] != INF && backup[u] + w < dist[v]) {
                    dist[v] = backup[u] + w;
                    relaxed = true;
                }
            }
            
            // 如果某跑完所有边都没有任何更新，说明已经触碰到了绝对最短路，可以直接中止
            if (!relaxed) break;
        }
    }

    long long get_dist(int target) {
        return dist[target];
    }
};

/*
int main() {
    int n = 3;
    BellmanFord graph(n);

    // 建图
    graph.add_edge(1, 2, 5);

    int start_node = 1;
    int k = 1; // 限制最多只能走 1 条边

    // 运行算法，传入起点和最大边数限制
    graph.run(start_node, k);

    long long ans = graph.get_dist(3);
    
    // 如果 k=1，1->3 只能走 15 的直达边，结果是 15。
    // 如果 k=2，1->2->3 才能走通，结果会变成 5+5=10。
    if (ans == graph.INF) {
        cout << "在 " << k << " 步内不可达\n";
    } else {
        cout << "在最多走 " << k << " 条边的情况下，最短距离为: " << ans << "\n";
    }
    return 0;
}
*/	
	
    \end{lstlisting}
    \subsubsection{求树的直径与找树心}
    \begin{lstlisting}
   const int MAXN = 100005;
vector<int> adj[MAXN];

int max_dist;
int farthest_node;
int parent_node[MAXN]; // 用于记录路径

// 第一次 DFS：只找最远点
void dfs1(int u, int p, int d) {
    if (d > max_dist) {
        max_dist = d;
        farthest_node = u;
    }
    for (int v : adj[u]) {
        if (v != p) {
            dfs1(v, u, d + 1);
        }
    }
}

// 第二次 DFS：找最远点，并记录父节点以便回溯路径
void dfs2(int u, int p, int d) {
    parent_node[u] = p; // 记录是从哪个节点走到 u 的
    if (d > max_dist) {
        max_dist = d;
        farthest_node = u;
    }
    for (int v : adj[u]) {
        if (v != p) {
            dfs2(v, u, d + 1);
        }
    }
}

void find_tree_centers() {
    // 1. 第一次 DFS 找直径的一个端点 A
    max_dist = -1;
    dfs1(1, 0, 0);
    int A = farthest_node;

    // 2. 第二次 DFS 找直径的另一个端点 B，并记录路径
    max_dist = -1;
    dfs2(A, 0, 0);
    int B = farthest_node;

    // 3. 从 B 往回走到 A，提取整条直径的路径
    vector<int> path;
    int curr = B;
    while (curr != 0) {
        path.push_back(curr);
        curr = parent_node[curr];
    }
   // 4. 根据路径长度（节点数）找树心
    /*int len = path.size();
    if (len % 2 == 1) {
        // 直径上有奇数个节点（即直径长度为偶数），树心只有 1 个
        cout << "树心是: " << path[len / 2] << "\n";
    } else {
        // 直径上有偶数个节点（即直径长度为奇数），树心有 2 个
        cout << "树心是: " << path[len / 2 - 1] << " 和 " << path[len / 2] << "\n";
    }*/
}

注意两次dfs求树的直径只能在不带负权边的情况下用，如果带负权边就用树形dp。上面的代码是边权为1的情况，边权不全为1就搞个node，记录一下to和w，把dfs里面+1改成+w即可。树形dp的优点是可以求负权，缺点就是只能求直径了，其他什么都干不了，记录直径上的点什么的都做不了。

//树形dp求直径
struct Edge {
    int to;
    long long weight;
};

vector<Edge> adj[MAXN];
long long dp[MAXN]; 
long long tree_diameter; 

void dfs_dp(int u, int p) {
    // 如果允许路径长度为 0（即一个点也不走），dp 初始值为 0。
    // 如果题目要求必须包含边且全是负权，这里可能需要初始化为 -INF
    dp[u] = 0; 
    
    for (auto& edge : adj[u]) {
        int v = edge.to;
        long long w = edge.weight;
        
        if (v == p) continue;
        
        // 1. 先递归算出子节点 v 的 dp 值
        dfs_dp(v, u);
        
        // 2. 拼合路径，尝试更新全局直径
        // 此时 dp[u] 代表的是之前遍历过的其他子树的最长链
        tree_diameter = max(tree_diameter, dp[u] + dp[v] + w);
        
        // 3. 维护 dp[u] 本身，供 u 的父节点使用
        dp[u] = max(dp[u], dp[v] + w);
    }
}

/*void solve() {
    int n;
    cin >> n;
    
    // 初始化清空
    for (int i = 1; i <= n; i++) {
        adj[i].clear();
    }
    
    for (int i = 0; i < n - 1; i++) {
        int u, v;
        long long w;
        cin >> u >> v >> w;
        adj[u].push_back({v, w});
        adj[v].push_back({u, w});
    }
    
    tree_diameter = 0; // 视题目情况，如果全负权可能需要初始化为很小的负数
    dfs_dp(1, 0);
    cout << tree_diameter << "\n";
}*/
    \end{lstlisting}
    \subsubsection{LCA}
    \begin{lstlisting}
    //LCA处理森林也是没问题的，加个vis数组，dfs的时候更新一下vis，然后不只是对根节点dfs了，要遍历所有节点，遇到!vis[i]就对他dfs一遍。然后就是求lca或者路径最大值什么的，先要判断一下两个点在不在一个连通块，这个也非常简单。如果是kru搞出来的mst上做lca，那kru有现成的并查集，判断mst.find(u)!=mst.find(v)就行，没有的话就搞一个id数组，在dfs的时候给每个点标记一下他们的id，就是属于哪个树，判断id[u]!=id[v]就行。
       struct LCA {
    int n, LOG;
    vector<vector<int>> adj;
    vector<int> depth;
    vector<vector<int>> fa;

    // 构造函数：初始化大小，动态计算需要的最大倍增层数
    LCA(int _n) {
        n = _n;
        LOG = std::__lg(n) + 2; // __lg 返回下取整的对数，+2 保证安全裕度
        adj.resize(n + 1);
        depth.assign(n + 1, 0);
        fa.assign(n + 1, vector<int>(LOG, 0));
    }               

    // 添边（无向树）
    void add_edge(int u, int v) {
        adj[u].push_back(v);
        adj[v].push_back(u);
    }

    // DFS 预处理深度和倍增表
    void dfs(int u, int p, int d) {
        depth[u] = d;
        fa[u][0] = p; 
        
        for (int k = 1; k < LOG; k++) {
            if (fa[u][k - 1] != 0) {
                fa[u][k] = fa[fa[u][k - 1]][k - 1];
            } else {
                fa[u][k] = 0; 
            }
        }
        
        for (int v : adj[u]) {
            if (v != p) {
                dfs(v, u, d + 1);
            }
        }
    }

    // 外部调用接口：建树初始化
    void build(int root = 1) {
        dfs(root, 0, 1);
    }

    // 查询 u 和 v 的最近公共祖先
    int get_lca(int u, int v) {
        if (depth[u] < depth[v]) swap(u, v);
        
        // 1. 深度对齐
        for (int k = LOG - 1; k >= 0; k--) {
            if (depth[u] - (1 << k) >= depth[v]) {
                u = fa[u][k];
            }
        }
        
        if (u == v) return u;
        
        // 2. 同步逼近
        for (int k = LOG - 1; k >= 0; k--) {
            if (fa[u][k] != fa[v][k]) {
                u = fa[u][k];
                v = fa[v][k];
            }
        }
        return fa[u][0];
    }

    // 附加实用功能：求树上任意两点的距离（边数）
    int get_dist(int u, int v) {
        return depth[u] + depth[v] - 2 * depth[get_lca(u, v)];
    }
};

/*
// 调用示例：
void solve() {
    int n, q; // n个节点，q次查询
    cin >> n >> q;
    
    LCA tree(n);
    
    for (int i = 0; i < n - 1; i++) {
        int u, v;
        cin >> u >> v;
        tree.add_edge(u, v);
    }
    
    tree.build(1); // 以 1 为根节点预处理
    
    while (q--) {
        int u, v;
        cin >> u >> v;
        cout << tree.get_lca(u, v) << "\n";
    }
}
*/
    \end{lstlisting}
    \subsubsection{Kruskal}
    \begin{lstlisting}
        struct Kru {
    struct Edge {
        int u, v;
        long long w;
        bool operator<(const Edge& other) const {
            return w < other.w;
        }
    };

    int n;
    vector<Edge> edges;       // 存储原图所有边
    vector<Edge> mst_edges;   // 存储构成最小生成树的边
    vector<int> fa;           // 并查集数组（一维 vector，绝对不会 MLE）
    long long min_weight;     // 最小生成树的总权值
    bool is_connected;        // 图是否连通

    // 构造函数：初始化节点个数和并查集
    Kru(int _n) {
        n = _n;
        fa.resize(n + 1);
        for (int i = 0; i <= n; i++) {
            fa[i] = i;
        }
        min_weight = 0;
        is_connected = false;
    }

    // 添加边
    void add_edge(int u, int v, long long w) {
        edges.push_back({u, v, w});
    }

    // 并查集带路径压缩的查找
    int find(int x) {
        return fa[x] == x ? x : fa[x] = find(fa[x]);
    }

    // 执行算法，返回最小生成树的总权值（如果不连通，返回 -1）
    long long build() {
        // 1. 将所有边按权值从小到大排序
        sort(edges.begin(), edges.end());
        
        mst_edges.clear();
        min_weight = 0;
        int edge_cnt = 0;

        // 2. 贪心选边
        for (const auto& e : edges) {
            int root_u = find(e.u);
            int root_v = find(e.v);
            
            if (root_u != root_v) {
                fa[root_u] = root_v;      // 合并连通块
                min_weight += e.w;        // 累加权值
                mst_edges.push_back(e);   // 记录这条边
                edge_cnt++;
                
                // 如果已经选了 n-1 条边，一棵树就建成了，提前退出
                if (edge_cnt == n - 1) break;
            }
        }
        
        // 3. 检查是否选够了 n-1 条边
        is_connected = (edge_cnt == n - 1 || n == 1);
        return is_connected ? min_weight : -1;
    }
};

/*
// 调用示例：
void solve() {
    int n, m; 
    cin >> n >> m;
    
    Kru mst(n);
    
    for(int i = 0; i < m; i++) {
        int u, v;
        long long w;
        cin >> u >> v >> w;
        mst.add_edge(u, v, w);
    }
    
    long long ans = mst.build();
    
    if (ans == -1) {
        cout << "orz" << "\n"; // 图不连通
    } else {
        cout << ans << "\n";
        
        // 如果题目需要你在 MST 上跑 LCA，可以这样直接用：
        // LCA tree(n);
        // for (auto& e : mst.mst_edges) {
        //     tree.add_edge(e.u, e.v); // 注意这里可能要根据情况传边权
        // }
        // tree.build();
    }
}
*/   
    \end{lstlisting}    
    \subsubsection{Prim}
    \begin{lstlisting}
    //注意prim的边加入的时候不是从小到大，是无序的，不像kruskal那样从小到大加边，因为prim是通过点去加。所以通过prim去做聚类问题的时候可以先把n-1条边给搞出来，开个edge数组记录一下，然后排下序，再取第k个，第n-k个，第k+1（根据题目来）个值什么的。edge数组记录的时候注意：prim的第一轮是会空转一轮先加个0进去的，所以加边的时候可以加个i!=1,把第一轮的0忽略掉。
    struct Prim {
    int n;
    vector<long long> dis; // dis[i] 表示点 i 距离当前生成树集合的最短距离
    vector<bool> vis;      // vis[i] 表示点 i 是否已经加入生成树
    long long min_weight;  // 最小生成树总权值
    bool is_connected;     // 图是否连通

    // 【针对隐式坐标图】这里用来存点的坐标，避免开 V^2 的二维矩阵
    // vector<pair<long long, long long>> points;

    Prim(int _n) {
        n = _n;
        dis.assign(n + 1, INF);
        vis.assign(n + 1, false);
        min_weight = 0;
        is_connected = false;
    }

    // 获取 u 和 v 之间的边权（核心拓展点）
    inline long long get_weight(int u, int v) {
        // 场景 A（纯坐标题）：不要存边，直接现算！
        //外面先开个vector<pair<int,int>>point来存点的坐标
        // long long dx = points[u].first - points[v].first;
        // long long dy = points[u].second - points[v].second;
        // return dx * dx + dy * dy; // 如果题目求的是欧几里得距离平方

        // 场景 B（普通的给定边稠密图）：
        // return G[u][v]
        //G是在外面开的二维数组，提前存每个点两两之间的边权
        return 0; 
    }

    // 执行算法，返回 MST 总权值，不连通返回 -1
    long long build(int start_node = 1) {
        min_weight = 0;
        dis[start_node] = 0;
        int cnt = 0; // 记录成功加入了多少个点

        for (int i = 1; i <= n; i++) {
            long long min_d = INF;
            int u = -1;

            // 1. 扫描一遍未加入的点，找离当前生成树最近的
            for (int j = 1; j <= n; j++) {
                if (!vis[j] && dis[j] < min_d) {
                    min_d = dis[j];
                    u = j;
                }
            }

            // 如果找不到，说明剩下的点和树断开了，图不连通
            if (u == -1) break;

            // 2. 将 u 吞并到生成树集合中
            vis[u] = true;
            min_weight += min_d;
            cnt++;

            // 3. 用刚加入的 u 去更新外面那些未加入点到树的距离
            for (int v = 1; v <= n; v++) {
                if (!vis[v]) {
                    long long w = get_weight(u, v); // 动态获取距离
                    if (w < dis[v]) {
                        dis[v] = w;
                    }
                }
            }
        }

        is_connected = (cnt == n);
        return is_connected ? min_weight : -1;
    }
};

/*
// 调用示例（以坐标建图为例）：
void solve() {
    int n; 
    cin >> n;
    
    Prim mst(n);
    // mst.points.resize(n + 1);
    // for(int i = 1; i <= n; i++) cin >> mst.points[i].first >> mst.points[i].second;
    
    long long ans = mst.build();
    
    if (ans == -1) cout << "orz\n";
    else cout << ans << "\n";
}
*/   
    \end{lstlisting}
    \subsubsection{严格次小生成树}
    \begin{lstlisting}
    struct Kru {
    struct Edge {
        int u, v;
        long long w;
        bool operator<(const Edge& other) const {
            return w < other.w;
        }
    };

    int n;
    vector<Edge> edges;       // 存储原图所有边
    vector<Edge> mst_edges;   // 存储构成最小生成树的边
    vector<Edge>nedges;
    vector<int> fa;           // 并查集数组（一维 vector，绝对不会 MLE）
    long long min_weight;     // 最小生成树的总权值
    bool is_connected;        // 图是否连通

    // 构造函数：初始化节点个数和并查集
    Kru(int _n) {
        n = _n;
        fa.resize(n + 1);
        for (int i = 0; i <= n; i++) {
            fa[i] = i;
        }
        min_weight = 0;
        is_connected = false;
    }

    // 添加边
    void add_edge(int u, int v, long long w) {
        edges.push_back({u, v, w});
    }

    // 并查集带路径压缩的查找
    int find(int x) {
        return fa[x] == x ? x : fa[x] = find(fa[x]);
    }

    // 执行算法，返回最小生成树的总权值（如果不连通，返回 -1）
    long long build() {
        // 1. 将所有边按权值从小到大排序
        sort(edges.begin(), edges.end());
        
        mst_edges.clear();
        min_weight = 0;
        int edge_cnt = 0;

        // 2. 贪心选边
        for (const auto& e : edges) {
            int root_u = find(e.u);
            int root_v = find(e.v);
            
            if (root_u != root_v) {
                fa[root_u] = root_v;      // 合并连通块
                min_weight += e.w;        // 累加权值
                mst_edges.push_back(e);   // 记录这条边
                edge_cnt++;
                
            }
            else{
                nedges.pb(e);
            }
        }
        return min_weight;
        // 3. 检查是否选够了 n-1 条边
        //is_connected = (edge_cnt == n - 1 || n == 1);
        //return is_connected ? min_weight : -1;
    }
};

/*
// 调用示例：
void solve() {
    int n, m; 
    cin >> n >> m;
    
    Kru mst(n);
    
    for(int i = 0; i < m; i++) {
        int u, v;
        long long w;
        cin >> u >> v >> w;
        mst.add_edge(u, v, w);
    }
    
    long long ans = mst.build();
    
    if (ans == -1) {
        cout << "orz" << "\n"; // 图不连通
    } else {
        cout << ans << "\n";
        
        // 如果题目需要你在 MST 上跑 LCA，可以这样直接用：
        // LCA tree(n);
        // for (auto& e : mst.mst_edges) {
        //     tree.add_edge(e.u, e.v); // 注意这里可能要根据情况传边权
        // }
        // tree.build();
    }
}
*/
		       struct LCA {
    int n, LOG;
    vector<vector<pii>> adj;
    vector<int> depth;
    vector<vector<int>> fa;
    vector<vector<int>>mx1,mx2;

    // 构造函数：初始化大小，动态计算需要的最大倍增层数
    LCA(int _n) {
        n = _n;
        LOG = std::__lg(n) + 2; // __lg 返回下取整的对数，+2 保证安全裕度
        adj.resize(n + 1);
        depth.assign(n + 1, 0);
        fa.assign(n + 1, vector<int>(LOG, 0));
        mx1.assign(n+1,vector<int>(LOG,-1));
        mx2.assign(n+1,vector<int>(LOG,-1));
    }               

    // 添边（无向树）
    void add_edge(int u, int v,int w) {
        adj[u].push_back({v,w});
        adj[v].push_back({u,w});
    }

    // DFS 预处理深度和倍增表
    void dfs(int u, int p, int d,int ww) {
        depth[u] = d;
        fa[u][0] = p; 
        mx1[u][0]=ww,mx2[u][0]=-1;
        for (int k = 1; k < LOG; k++) {
            if (fa[u][k - 1] != 0) {
                fa[u][k] = fa[fa[u][k - 1]][k - 1];
                int a1=mx1[u][k-1],a2=mx2[u][k-1];
                int b1=mx1[fa[u][k-1]][k-1],b2=mx2[fa[u][k-1]][k-1];
                mx1[u][k]=max(a1,b1);
                if(a1<b1){
                    mx2[u][k]=max(a1,b2);
                }
                else if(a1>b1){
                    mx2[u][k]=max(a2,b1);
                }
                else{
                    mx2[u][k]=max(a2,b2);
                }
            } else {
                fa[u][k] = 0; 
            }
        }
        
        for (auto v : adj[u]) {
            if (v.first != p) {
                dfs(v.first, u, d + 1,v.second);
            }
        }
    }

    // 外部调用接口：建树初始化
    void build(int root = 1) {
        dfs(root, 0, 1,-1);
    }

    // 查询 u 和 v 的最近公共祖先
    pii get_res(int u, int v) {
         int res1=-1,res2=-1;
        auto merge=[&](int m1,int m2){
            if(m1>res1){
                res2=max(res1,m2);
                res1=m1;
            }
            else if(m1==res1){
                res2=max(res2,m2);
            }
            else if(m1<res1&&m1>res2){
                res2=m1;
            }
        };
        if (depth[u] < depth[v]) swap(u, v);
        // 1. 深度对齐
        for (int k = LOG - 1; k >= 0; k--) {
            if (depth[u] - (1 << k) >= depth[v]) {
                merge(mx1[u][k],mx2[u][k]);
                u = fa[u][k];
            }
        }
        
        if (u == v) return {res1,res2};
        
        // 2. 同步逼近
        for (int k = LOG - 1; k >= 0; k--) {
            if (fa[u][k] != fa[v][k]) {
                merge(mx1[u][k],mx2[u][k]);
                merge(mx1[v][k],mx2[v][k]);
                u = fa[u][k];
                v = fa[v][k];
            }
        }
        merge(mx1[u][0],mx2[u][0]);
        merge(mx1[v][0],mx2[v][0]);
        return {res1,res2};
    }

    // 附加实用功能：求树上任意两点的距离（边数）
    //int get_dist(int u, int v) {
        //return depth[u] + depth[v] - 2 * depth[get_lca(u, v)];
    //}
};

/*
// 调用示例：
void solve() {
    int n, q; // n个节点，q次查询
    cin >> n >> q;
    
    LCA tree(n);
    
    for (int i = 0; i < n - 1; i++) {
        int u, v;
        cin >> u >> v;
        tree.add_edge(u, v);
    }
    
    tree.build(1); // 以 1 为根节点预处理
    
    while (q--) {
        int u, v;
        cin >> u >> v;
        cout << tree.get_lca(u, v) << "\n";
    }
}
*/
		void solve()
		{		
			int n,m;
            cin>>n>>m;
            Kru mst(n);
            for(int i=0;i<m;i++){
                int u,v,w;
                cin>>u>>v>>w;
                if(u==v) continue;
                mst.add_edge(u,v,w);
            }
            int ini=mst.build();
            LCA tr(n);
            for(auto e:mst.mst_edges){
                tr.add_edge(e.u,e.v,e.w);
            }
            tr.build();
            int ans=2e18;
            for(auto e:mst.nedges){
                int u=e.u;
                int v=e.v;
                int w=e.w;
                int r1=tr.get_res(u,v).first;
                int r2=tr.get_res(u,v).second;
                if(w>r1) ans=min(ans,ini+w-r1);
                else if(r2!=-1) ans=min(ans,ini+w-r2);
            }
            cout<<ans<<"\n";
		}
        如果是非严格次小生成树，那直接维护一个最大值就够了，更简单
    \end{lstlisting}
    \subsubsection{dagdp求dag上的最长链以及有几条最长链}
    \begin{lstlisting}
        vector<int>a=ts.order;
			int l=a.size();
			vector<int>dp(l+1,1);
			vector<int>cnt(l+1,1);
			for(int i=1;i<=l;i++){
				dp[i]=g.scc_sz[i];
			}
			for(int i=0;i<l;i++){
				int u=a[i];
				for(auto x:ts.adj[u]){
					if(dp[u]+g.scc_sz[x]>dp[x]){
						dp[x]=dp[u]+g.scc_sz[x];//这里是tarjan缩点之后每个超级点带了一个权值的情况，如果正常dag就++就行了。
						cnt[x]=cnt[u];
					}
					else if(dp[u]+g.scc_sz[x]==dp[x]){
						cnt[x]=(cnt[x]+cnt[u])%w;
					}
				}
			}
			int mx=0;
			for(int i=1;i<=l;i++) mx=max(mx,dp[i]);
			int ans=0;
			for(int i=1;i<=l;i++){
				if(dp[i]==mx){
					ans=(ans+cnt[i])%w;
				}
			}
			cout<<mx<<"\n";//最长链的节点数
			cout<<ans<<"\n";//有几条最长链
    \end{lstlisting}
    \subsubsection{scc}
    \begin{lstlisting}
        struct scc {
    int n; // 节点数
    vector<vector<int>> adj; // 有向图邻接表
    vector<int> dfn, low, scc_id, scc_sz; // Tarjan 核心数组
    vector<bool> in_stk; // 标记是否在栈中
    stack<int> stk;
    int timer, scc_cnt; // 时间戳和强连通分量数量

    // 构造函数，初始化大小
    scc(int _n) {
        n = _n;
        adj.resize(n + 1);
        dfn.assign(n + 1, 0);
        low.assign(n + 1, 0);
        scc_id.assign(n + 1, 0);  
        scc_sz.assign(n + 1, 0);
        in_stk.assign(n + 1, false);
        timer = 0;
        scc_cnt = 0;
    }

    // 添加有向边
    void add_edge(int u, int v) {
        adj[u].push_back(v);
    }

    // Tarjan 核心 DFS 过程（内部调用）
    void dfs(int u) {
        dfn[u] = low[u] = ++timer;
        stk.push(u);
        in_stk[u] = true;

        for (int v : adj[u]) {
            if (!dfn[v]) {
                dfs(v);
                low[u] = min(low[u], low[v]);
            } else if (in_stk[v]) {
                low[u] = min(low[u], dfn[v]);
            }
        }

        if (dfn[u] == low[u]) {
            ++scc_cnt;
            int v;
            do {
                v = stk.top();
                stk.pop();
                in_stk[v] = false;
                scc_id[v] = scc_cnt; // 记录原节点属于哪个超级节点
                scc_sz[scc_cnt]++;   // 记录该超级节点包含几个原节点
            } while (v != u);
        }
    }

    // 运行 Tarjan 算法求出所有强连通分量
    void run() {
        for (int i = 1; i <= n; i++) {
            if (!dfn[i]) {
                dfs(i);
            }
        }
    }
};

/*
int main() {
    int n, m;
    cin >> n >> m;

    // 实例化结构体，传入节点数量
    SCC graph(n);

    // 读入有向边
    for (int i = 0; i < m; ++i) {
        int u, v;
        cin >> u >> v;
        graph.add_edge(u, v); 
    }

    // 运行算法
    graph.run();

    // 获取结果示例
    cout << "整张图被缩成了 " << graph.scc_cnt << " 个超级节点\n";
    
    // 如果想要知道两个点是否互相可达（即在同一个强连通分量里）
    int u = 1, v = 2;
    if (graph.scc_id[u] == graph.scc_id[v]) {
        cout << u << " 和 " << v << " 在同一个强连通分量中\n";
    }

    return 0;
}
*/
    \end{lstlisting}
    \subsubsection{bc(bridge\&cut 求桥和割点)}
    \begin{lstlisting}
        struct bc {
    int n, timer, ecnt;
    vector<vector<pair<int,int>>> adj;

    vector<int> dfn, low;
    vector<int> br,cut;

    bc(int _n) {
        n = _n;
        timer = ecnt = 0;
        adj.resize(n + 1);
    }

    // 添加无向边
    void add_edge(int u, int v) {
        int id = ecnt++;
        adj[u].push_back({v, id});
        adj[v].push_back({u, id});
    }

    void dfs(int u, int pe) {
        dfn[u] = low[u] = ++timer;

        int child = 0;

        for (auto [v, id] : adj[u]) {
            if (id == pe) continue;

            if (!dfn[v]) {
                child++;

                dfs(v, id);

                low[u] = min(low[u], low[v]);

                // 判断桥
                if (low[v] > dfn[u]) {
                    br[id] = 1;
                }

                // 判断非根节点割点
                if (pe != -1 && low[v] >= dfn[u]) {
                    cut[u] = 1;
                }
            }
            else {
                low[u] = min(low[u], dfn[v]);
            }
        }

        // 根节点特殊判断
        if (pe == -1 && child >= 2) {
            cut[u] = 1;
        }
    }

    void build() {
        timer = 0;

        dfn.assign(n + 1, 0);
        low.assign(n + 1, 0);

        br.assign(ecnt, 0);
        cut.assign(n + 1, 0);

        // 处理非连通图
        for (int i = 1; i <= n; i++) {
            if (!dfn[i]) {
                dfs(i, -1);
            }
        }
    }
};
/*用法示例
bc g(n);
g.add_edge(u, v);
g.build();

g.br[id]; // 是否为桥
g.cut[u]; // 是否为割点*/
    \end{lstlisting}
    \subsubsection{ap(割点)}
    \begin{lstlisting}
        struct ap {
    int n;
    vector<vector<int>> adj;
    vector<int> dfn, low;
    vector<bool> is_cut;
    int timer;
    int cut_cnt; // 割点总数

    ap(int _n) {
         n = _n;
        adj.assign(n + 1, vector<int>());
        dfn.assign(n + 1, 0);
        low.assign(n + 1, 0);
        is_cut.assign(n + 1, false);
        timer = 0;
        cut_cnt = 0;
    }

    void add_edge(int u, int v) {
        adj[u].push_back(v);
        adj[v].push_back(u);
    }

    // p 为父节点，用于防止顺着无向边往回走
    void dfs(int u, int p) {
        dfn[u] = low[u] = ++timer;
        int child_cnt = 0; // 记录搜索树上的独立子节点数量

        for (int v : adj[u]) {
            if (v == p) continue; // 绝不能用父节点的 dfn 更新自己的 low

            if (!dfn[v]) { // 分支1：树枝边（v 没去过）
                child_cnt++;
                dfs(v, u);
                low[u] = min(low[u], low[v]); // 用小弟的 low 更新自己

                // 核心判定 1：非根节点，且小弟无论怎么绕也跨不过我
                if (p != 0 && low[v] >= dfn[u]) {
                    is_cut[u] = true;
                }
            } else { // 分支2：返祖边（v 去过，且不是父节点）
                low[u] = min(low[u], dfn[v]); // 只能用 dfn 更新，不能用 low
            }
        }

        // 核心判定 2：如果是 DFS 树的根节点，且有两个以上的独立分支
        if (p == 0 && child_cnt >= 2) {
            is_cut[u] = true;
        }
    }

    // 图可能不连通，需要遍历所有节点
    void build() {
        for (int i = 1; i <= n; i++) {
            if (!dfn[i]) {
                dfs(i, 0); // 0 代表没有父节点，即作为根节点出发
            }
        }
        
        // 统计割点总数
        for (int i = 1; i <= n; i++) {
            if (is_cut[i]) cut_cnt++;
        }
    }
};

/*
void solve() {
    int n, m;
    cin >> n >> m;
    
    CutVertex graph(n);
    
    for (int i = 0; i < m; i++) {
        int u, v;
        cin >> u >> v;
        // 遇到重边不用怕，这套逻辑天然免疫重边干扰
        graph.add_edge(u, v);
    }
    
    graph.build();
    
    cout << graph.cut_cnt << "\n";
    for (int i = 1; i <= n; i++) {
        if (graph.is_cut[i]) {
            cout << i << " ";
        }
    }
    cout << "\n";
}
*/
    \end{lstlisting}
    \subsubsection{vbcc(点双)}
    \begin{lstlisting}
        struct vbcc {
    int n; // 节点数
    vector<vector<int>> adj; // 邻接表
    vector<int> dfn, low, st; // Tarjan 数组与手写栈
    int timer;
    
    vector<bool> is_cut; // 标记某个点是否为全局割点
    vector<vector<int>> bccs; // 存储所有的点双连通分量

    // 构造函数，初始化大小
    vbcc(int _n) {
        n = _n;
        adj.resize(n + 1);
        dfn.assign(n + 1, 0);
        low.assign(n + 1, 0);
        is_cut.assign(n + 1, false);
        timer = 0;
    }

    // 添加无向边
    void add_edge(int u, int v) {
        adj[u].push_back(v);
        adj[v].push_back(u);
    }

    // Tarjan 核心逻辑
    void tarjan(int u, int root) {
        dfn[u] = low[u] = ++timer;
        st.push_back(u);
        int child = 0;

        // 特判孤立点：没有连边的点自身构成一个点双
        if (u == root && adj[u].empty()) {
            bccs.push_back({u});
            return;
        }

        for (int v : adj[u]) {
            if (!dfn[v]) {
                child++;
                tarjan(v, root);
                low[u] = min(low[u], low[v]);

                // 发现点双边界
                if (low[v] >= dfn[u]) {
                    is_cut[u] = true; // 暂时标记，若是根节点稍后特判修正
                    vector<int> comp;
                    
                    // 弹栈直到弹出 v 为止
                    while (true) {
                        int x = st.back();
                        st.pop_back();
                        comp.push_back(x);
                        if (x == v) break;
                    }
                    // u 作为割点也要加入该点双，但 u 本身不出栈
                    comp.push_back(u);
                    bccs.push_back(comp);
                }
            } else {
                // 返祖边更新 low 值
                low[u] = min(low[u], dfn[v]);
            }
        }

        // 特判根节点是否真的是割点（至少有两个独立的子树）
        if (u == root && child < 2) {
            is_cut[u] = false;
        }
    }

    // 运行 Tarjan 算法求解图中的所有 v-BCC 和割点
    void run() {
        for (int i = 1; i <= n; ++i) {
            if (!dfn[i]) {
                // 每次未访问的点都是一个新的连通块的 root
                tarjan(i, i);
            }
        }
    }
};

/*
    int n, m;
    cin >> n >> m;

    // 1. 实例化结构体，传入节点数量
    VBCC graph(n);

    // 2. 读入无向边
    for (int i = 0; i < m; ++i) {
        int u, v;
        cin >> u >> v;
        graph.add_edge(u, v);
    }

    // 3. 运行算法
    graph.run();

    // 4. 获取结果：输出图中有多少个割点
    int cut_count = 0;
    for (int i = 1; i <= n; ++i) {
        if (graph.is_cut[i]) {
            cut_count++;
        }
    }
    cout << "割点总数: " << cut_count << "\n";

    // 5. 获取结果：输出所有的点双连通分量
    cout << "点双连通分量总数: " << graph.bccs.size() << "\n";
    for (int i = 0; i < graph.bccs.size(); ++i) {
        cout << "BCC #" << i + 1 << ": ";
        for (int node : graph.bccs[i]) {
            cout << node << " ";
        }
        cout << "\n";
    }

//圆方树
    int total_nodes = n + g.bccs.size(); 
    //完整圆方树
    vector<vector<int>> tree(total_nodes + 1); // 邻接表开到最大节点数
    for (int i = 0; i < g.bccs.size(); i++) {
        int square = n + i + 1; // 方点编号从 n+1 开始
        
        for (int round : g.bccs[i]) {
            tree[square].push_back(round);
            tree[round].push_back(square);
        }
    }
    //压缩圆方树
    vector<vector<int>> comp_tree(total_nodes + 1); 
    for (int i = 0; i < g.bccs.size(); i++) {
        int square = n + i + 1; 
        
        for (int round : g.bccs[i]) {
            // 方点【仅】向该点双内的【割点】连边
            if (g.is_cut[round]) {
                comp_tree[square].push_back(round);
                comp_tree[round].push_back(square);
            }
        }
    }
*/
    \end{lstlisting}
    \subsubsection{2-sat}
    \begin{lstlisting}
        //注意2sat add_edge的时候一定是成对出现的，因为你加了u->v，就一定要加v反->u反。一个命题成立那么它的逆否命题必然成立。
        struct tsat {
    int n; // 变量数量
    vector<vector<int>> adj; // 蕴含图

    // Tarjan 所需数组
    vector<int> dfn, low, scc_id, scc_sz;
    vector<bool> in_stk;
    stack<int> stk;
    int timer, scc_cnt;

    vector<bool> ans; // 最终答案，ans[i] 表示 xi 是否为 true

    // 构造函数
    tsat(int _n) {
        n = _n;
        adj.resize(2 * n + 1);

        dfn.assign(2 * n + 1, 0);
        low.assign(2 * n + 1, 0);
        scc_id.assign(2 * n + 1, 0);
        scc_sz.assign(2 * n + 1, 0);
        in_stk.assign(2 * n + 1, false);

        timer = 0;
        scc_cnt = 0;

        ans.assign(n + 1, false);
    }

    // 将逻辑文字转换为图中的节点
    // 正数 i 表示 xi 为 true
    // 负数 -i 表示 xi 为 false
    int node(int x) {
        if (x > 0) return x;
        return n + (-x);
    }

    // 求一个节点对应的反状态节点
    int neg(int u) {
        if (u <= n) return u + n;
        return u - n;
    }

    // 添加一条蕴含边：a -> b
    // a、b 均为带正负号的逻辑文字
    void add_edge(int a, int b) {
        int u = node(a);
        int v = node(b);
        adj[u].push_back(v);
    }

    // 添加一个或条件：(a | b)
    // 对应两条边：!a -> b，!b -> a
    void add_or(int a, int b) {
        int u = node(a);
        int v = node(b);

        adj[neg(u)].push_back(v);
        adj[neg(v)].push_back(u);
    }

    // 添加一个异或条件：(a ^ b)
    void add_xor(int a, int b) {
        add_or(a, b);
        add_or(-a, -b);
    }

    // 强制某个逻辑文字为 true
    void force(int a) {
        add_or(a, a);
    }

    // Tarjan 核心 DFS
    void dfs(int u) {
        dfn[u] = low[u] = ++timer;
        stk.push(u);
        in_stk[u] = true;

        for (int v : adj[u]) {
            if (!dfn[v]) {
                dfs(v);
                low[u] = min(low[u], low[v]);
            }
            else if (in_stk[v]) {
                low[u] = min(low[u], dfn[v]);
            }
        }

        if (dfn[u] == low[u]) {
            ++scc_cnt;
            int v;

            do {
                v = stk.top();
                stk.pop();

                in_stk[v] = false;
                scc_id[v] = scc_cnt;
                scc_sz[scc_cnt]++;
            } while (v != u);
        }
    }

    // 运行 Tarjan
    void run() {
        timer = 0;
        scc_cnt = 0;

        fill(dfn.begin(), dfn.end(), 0);
        fill(low.begin(), low.end(), 0);
        fill(scc_id.begin(), scc_id.end(), 0);
        fill(scc_sz.begin(), scc_sz.end(), 0);
        fill(in_stk.begin(), in_stk.end(), false);

        for (int i = 1; i <= 2 * n; i++) {
            if (!dfn[i]) {
                dfs(i);
            }
        }
    }
};
    // 判断是否有解，并恢复一组合法赋值
    // 返回 true 表示有解，false 表示无解
   /* bool solve() {
        run();

        for (int i = 1; i <= n; i++) {
            // 正反状态处于同一个 SCC，无解
            if (scc_id[i] == scc_id[i + n]) {
                return false;
            }
        }

        // 恢复具体赋值
        // 当前 Tarjan 编号方式：弹栈时递增
        for (int i = 1; i <= n; i++) {
            ans[i] = (scc_id[i] < scc_id[i + n]);
        }

        return true;
    }
*/

    \end{lstlisting}
    \section{语法注意事项}
    \subsection{函数里有默认参数}
    \begin{lstlisting}
    函数里有默认参数的时候要把默认参数堆到右边，因为如果放左边那你所有参数都得设个默认值反正挺麻烦的，所以统一放右边就好了，然后传参的时候传非默认参数就行。
    例如：
    double build(int tar,int start_node = 1) 
    mst.build(1);//这里的1是tar
    \end{lstlisting}
    \subsection{想用几个char拼接成一个字符串}
    \begin{lstlisting}
        如果直接写string=char1+char2+char3这样子是不行的，因为这样会把三个char的ascii码相加，得到的是整数而不是字符串。
        正确做法是，如果char是从一个string里面取的，可以用substr，s.substr(2,3)代表从s的下标2开始截取长度为3的字串，然后写str=str1+str2+str3...字符串相加是可以正常得到字符串的。或者这样写，string+=char1,string+=char2，这样是可以的。
    \end{lstlisting}
\end{document}

